\section{Soft Actor-Critic (SAC)}
\label{sec:sac}

\subsection{Algorithm Overview}

Soft Actor-Critic (SAC)~\cite{haarnoja2018soft} is an off-policy actor-critic algorithm that operates within the maximum entropy reinforcement learning framework. Unlike standard RL algorithms that solely aim to maximize the expected sum of rewards, SAC maximizes a dual objective consisting of both the expected return and the expected entropy of the policy:
\begin{equation}
    J(\pi) = \sum_{t=0}^{T} \mathbb{E}_{(s_t, a_t) \sim \rho_\pi} \left[ r(s_t, a_t) + \alpha \mathcal{H}\left(\pi(\cdot|s_t)\right) \right].
\end{equation}
Here, $\alpha$ is the temperature parameter that explicitly dictates the trade-off between exploitation (maximizing the reward) and exploration (maximizing the entropy $\mathcal{H}$). 

\begin{figure*}[htbp]
    \centering
    \includegraphics[width=\linewidth]{sac_hopper_reward_speed_x.pdf}
    \caption{Training dynamics of SAC on the Hopper environment (smoothed with an exponential moving average)}
    \label{fig:sac_results}
\end{figure*}

To optimize this objective, SAC concurrently learns a stochastic policy $\pi_\theta(a|s)$ and two parameterized soft Q-functions, $Q_{\phi_1}(s, a)$ and $Q_{\phi_2}(s, a)$. The use of twin Q-networks mitigates the overestimation bias commonly found in value-based methods, taking the minimum of the two Q-values during policy evaluation. The policy itself is typically parameterized as a squashed Gaussian distribution, ensuring that sampled actions inherently respect the bounded continuous action space. 

By aggressively reusing past experiences from a replay buffer, SAC theoretically bridges the gap between the high sample efficiency of off-policy methods (such as Rainbow) and the continuous control stability of on-policy methods (such as PPO). Furthermore, modern implementations of SAC automatically tune the temperature parameter $\alpha$ via dual optimization, ensuring that the exploration bounds dynamically adjust as the policy converges.

\subsection{Hopper: Learning from Pixels}
\label{sec:sac_hopper_results}

We evaluated SAC on the Hopper environment to directly compare its sample efficiency and computational demands against PPO and our value-based baselines. 

\textbf{Configuration:}
We maintain the foundational environment setup: $64 \times 64$ pixel inputs, a frame stack of 4, and the widened healthy\_angle\_range of $[-0.4, 0.4]$. Similar to our approach with PPO, we relied on Optuna to navigate the notoriously brittle hyperparameter landscape of continuous actor-critic methods. 
The fixed components include the replay buffer capacity and the network backbone configuration. Optuna tuned the following training parameters: actor\_lr = $10^{-5}$, critic\_lr = $7.8 \times 10^{-5}$, alpha\_lr = $4.5 \times 10^{-5}$, gamma = $0.99$, tau = $0.009$, batch\_size = $256$, gradient\_steps = $4$, start\_train\_after = $10000$, and share\_encoder = False.

\textbf{Results and Discussion:}
In terms of pure environment interactions, SAC is by far the most sample-efficient algorithm evaluated in this study. The optimized configuration successfully learns a highly robust jumping gait, reaching approximately 5 times the cumulative reward and final forward distance in just 500k frames. This drastically outpaces PPO (which required roughly 1.5M steps) and entirely eclipses standard DQN.

However, this exceptional sample efficiency comes at a severe computational cost. Because SAC performs computationally heavy gradient updates on three separate networks (the actor and two critics) at every environment step, its wall-clock training time is extremely slow, about 50x slower than PPO on the same number of environment frames.
Additionally, much like PPO, SAC exhibited high instability during the tuning phase. It is particularly sensitive to the initial temperature parameter and the actor's learning rate, frequently resulting in catastrophic policy divergence if the entropy collapses too early.

\subsection{Humanoid: Scaling to High-Dimensional Control}
\label{sec:sac_humanoid_results}

Given SAC's theoretical supremacy in sample-efficient continuous control, we attempted to scale it to the heavily constrained $96 \times 96$ pixel Humanoid environment. 

\textbf{Configuration:}
To give SAC the highest probability of success in this highly unstable bipedal domain, we relaxed the constraints on our Optuna search space. Recognizing that actor-critic methods are exceptionally sensitive to the scale and variance of the reward signal, we allowed Optuna to dynamically tune the reward weights ($w_{\text{vel}}$, $w_{\text{healthy}}$, and $w_{\text{ctrl}}$) alongside the core algorithmic parameters. 
Tuning reward weights proved critical for SAC's performance, improving learning on the humanoid environment. However, the computational demands of off-policy actor-critic methods on high-dimensional pixel inputs made it impossible to conduct extended training runs. We were limited to short-horizon trials (100k environment steps) to rapidly prune diverging configurations, which is insufficient for fully converging a stable movement policy in this complex environment.

\textbf{Results and Discussion:}
The inclusion of reward shaping in the hyperparameter search yielded a highly promising configuration. During our short-horizon Optuna trials, we observed a specific run that successfully managed to balance the humanoid and take a couple of stable, coordinated steps forward. This indicates that SAC is fundamentally capable of extracting a viable bipedal locomotion policy directly from high-dimensional visual features.

Unfortunately, we were unable to train this promising configuration to full convergence. The sheer computational overhead of off-policy pixel-based SAC on a 17-joint continuous action space is paralyzing. The time required to process the convolutional feature extraction, update the twin critics, and optimize the Gaussian actor at every single step made extended training horizons practically unfeasible within our compute budget. Consequently, while SAC demonstrates the highest theoretical ceiling for pixel-based continuous control, its deployment is severely bottlenecked by its immense wall-clock training time.